\chapter{Preprocessing}
\label{ch:preprocessing}

This chapter describes how the raw MetroAT recordings became the units of analysis, and
the design choices along the way that affect every later result.

\section{The dataset}

MetroAT~\cite{MetroAT:2026} is delivered as Parquet, so no conversion from CSV was needed.
The schema has 109 columns: 74 continuous sensor channels, 20 binary channels, 7
operational-mode signals, 4 failure and maintenance markers, and 4 timestamp columns. The
recordings run at 1\,Hz over roughly one year.

Several properties of the data shaped the processing:

\begin{itemize}
  \item \textbf{Normalised values.} All sensor channels are min-max normalised to
        approximately $[0,1]$; they are not physical units. Every threshold used in this
        study is therefore empirical or quantile-based, and no number in this report should
        be read as bar, \si{\metre\per\second\squared}, or \si{\kilo\pascal}.
  \item \textbf{Unsorted rows.} Rows are not guaranteed to be time-sorted within a daily
        file, so each file is sorted on load before any derivative is computed.
  \item \textbf{Streaming.} The data is processed one daily file at a time, keeping peak
        memory below 4\,GB rather than loading the full year at once.
  \item \textbf{Per-wagon replication.} Many pneumatic channels are duplicated across
        wagons, producing 214 sensor pairs with $|r| > 0.98$. No columns are deleted by
        hand. Instead, where a compact input is needed, principal component analysis is
        applied to the pneumatic set and the leading components are used in place of the raw
        channels; because so many channels are near-duplicates, a handful of components
        already captures most of the variance (Figure~\ref{fig:pca}).
\end{itemize}

\section{Train and held-out split}
\label{sec:split}

Models are fitted on a training window (June 2024 to 11 February 2025, about 14.3~million
rows) and scored on a held-out window from 12 February to June 2025 that the models never
see during fitting. Every out-of-sample figure reported later therefore comes from the
held-out window. The training window contains 10 documented failures and 11 maintenance
events, 8 of the failures being brake-system failures; these small counts are what make the
pre-failure analysis of Chapter~\ref{ch:failure} exploratory rather than conclusive.

\section{Windowing}
\label{sec:windowing}

The unit of analysis is a non-overlapping ten-second window. The roughly 30~million raw
rows reduce to about 1.49~million windows. For each window, summary statistics are computed
per channel (means and standard deviations, and minima and maxima for the kinematic
channels), so that each window carries a compact description of both motion and
air-system behaviour over its ten seconds. The ten-second choice is a compromise between
resolution and stability; its robustness is examined directly in
Section~\ref{sec:window-sensitivity}.

\section{Derived kinematics}

Train speed is the only raw motion channel. From it, acceleration is derived as the rate
of change of speed and jerk as the rate of change of acceleration. These derived channels,
summarised per window, are the sole basis for the operational states of
Chapter~\ref{ch:states}. Keeping the state definition to motion alone is the concrete
consequence of the feature-role split introduced in Section~\ref{sec:roles}.

\section{Feature roles}
\label{sec:roles}

This project runs clustering, classification, and regression over the \emph{same} set of
MetroAT columns: the operational states of Chapter~\ref{ch:states} are clustered from the
data, and later chapters predict those states and their braking intensity from it again.
Whenever the same columns can both define a target and be offered as predictors, leakage
becomes a real risk: a predictor that helped define the very thing it is asked to predict
will score near-perfectly and reveal nothing. To avoid this, every feature is assigned a
\emph{role} in a ledger (\texttt{feature\_roles.csv}) that fixes whether it may serve as a
predictor. The columns are split as follows:

\begin{itemize}
  \item \textbf{velocity}: the motion channels (speed, acceleration, jerk, and the derived
        deceleration-energy axis). These \emph{define} the operational states and the
        braking targets, so they are used only as targets and never as predictors.
  \item \textbf{actuation}: the brake command and its immediate pneumatic response
        (brake-cylinder pressure, proportional-valve pressure, spring-brake pressure,
        pneumatic braking force). Predictor-eligible, but closely tied to the brake demand.
  \item \textbf{auxiliary}: independent health channels that do not carry the brake command
        directly (main-reservoir pressure, load pressure, load signal,
        energy-braking-resistance, ambient temperature, compressor-running).
        Predictor-eligible, and the cleanest source of independent signal.
  \item \textbf{other}: timestamps, operating-mode flags, and failure/maintenance markers,
        used for the train/held-out split, time alignment, and validation, but never as
        predictors.
\end{itemize}

Table~\ref{tab:roles} summarises the ledger by role. The distinction between
\emph{actuation} and \emph{auxiliary} matters most in Chapter~\ref{ch:classification},
where the two are tested as separate predictor tiers: because the auxiliary channels never
touched the brake command, their ability to recover the deceleration state is the cleanest
test that the air system carries genuinely independent signal. The complete assignment of
channels to roles is listed in Appendix~\ref{app:roles}.

\begin{table}[htbp]
  \centering
  \caption{Sensor roles from the role ledger, aggregated across the window-level and
  event-level feature sets. Only \emph{actuation} and \emph{auxiliary} features are
  predictor-eligible; \emph{velocity} features define the targets and are never used as
  predictors.}
  \label{tab:roles}
  \begin{tabular}{lll}
    \toprule
    Role       & Predictor-eligible & Example channels \\
    \midrule
    velocity   & no  & train speed, acceleration, jerk, deceleration energy \\
    actuation  & yes & brake-cylinder / proportional-valve / spring-brake pressure, \\
               &     & pneumatic braking force \\
    auxiliary  & yes & main-reservoir pressure, load pressure, load signal, \\
               &     & energy-braking-resistance, ambient temperature, compressor \\
    other      & no  & timestamps, operating-mode flags, failure / maintenance markers \\
    \bottomrule
  \end{tabular}
\end{table}

\begin{figure}[htbp]
  \centering
  \includegraphics[width=0.8\linewidth]{pca_explained_variance.png}
  \caption{Cumulative explained variance of the pneumatic sensor set under PCA. Because 214
  sensor pairs are near-duplicates ($|r| > 0.98$) from per-wagon replication, a small number
  of components captures most of the variation, which is why the redundancy can be absorbed
  by PCA without dropping any columns.}
  \label{fig:pca}
\end{figure}